% ===========================================================================
\chapter{Experimental Results and Analysis}
% ===========================================================================

\section{Experimental Setup}

\subsection{Hardware and Software}

Experiments were conducted on a workstation running Windows 10 with Python 3.12. Table~\ref{tab:eval_summary} summarizes mean evaluation outcomes for Maskable PPO and the Best-Fit baseline under the setup described in this section.

\begin{table}[!ht]
\centering
\caption{Comparative evaluation summary (Best-Fit vs.\ Maskable PPO)}
\label{tab:eval_summary}
\begin{tabularx}{\textwidth}{X r r r l}
\hline
\textbf{Metric} & \textbf{Best-Fit mean} & \textbf{PPO mean} & \textbf{$\Delta$\%} & \textbf{Winner} \\ \hline
Acceptance Ratio & 0.9370 & 0.9236 & $-1.43$\% & Best-Fit \\
Avg Security Cost Heuristic & 0.0012 & 0.0010 & $-11.09$\% & PPO \\
Avg Security Margin & 1.0850 & 1.6257 & $+49.84$\% & PPO \\
Total Revenue & 9072.20 & 8965.40 & $-1.18$\% & Best-Fit \\
Avg Revenue / Request & 9.0722 & 8.9654 & $-1.18$\% & Best-Fit \\
Avg SFC Tenancy (SFCs/occupied node) & 1.0226 & 2.2189 & $+116.98$\% & Best-Fit \\
Avg VNF Tenancy (VNFs/occupied node) & 3.0672 & 3.8079 & $+24.15$\% & Best-Fit \\

\hline
\end{tabularx}
\end{table}

\subsection{Network Configuration}

The substrate network consists of 25 nodes connected as a complete graph ($p = 1.0$) with resources sampled uniformly as described in Section~3.1.3. SFC requests consist of chains of 3 VNFs with demands and constraints sampled as described in Section~3.2. The TTL follows a truncated exponential distribution on $[1, 25]$ with decay mean at the midpoint.

\subsection{Evaluation Protocol}

Evaluation uses 1000 SFC requests on the training substrate with a fresh request generator. To ensure fair comparison between the RL agent and the Best-Fit baseline, the Python and NumPy random number generator states are snapshotted before each algorithm's evaluation and reset to the same initial state, ensuring both algorithms process identical request sequences. The RL agent is evaluated in a single long episode (1000 requests, no resets), while the baseline uses a sustained-load simulation with the same substrate.

\section{Training Dynamics}

\subsection{Acceptance Ratio Progression}

Figure~\ref{fig:ar} illustrates the training acceptance ratio (rolling average over episodes) plotted against the number of training timesteps. An upward trend confirms that the agent is learning an increasingly effective placement policy over the course of training.

\begin{figure}[!ht]
\centering
\includegraphics[width=0.8\textwidth]{models/sfc_ppo_acceptance_ratio.png}
\caption{Rolling acceptance ratio during training. The acceptance ratio improves as the agent refines its placement strategy.}
\label{fig:ar}
\end{figure}

\section{Comparative Evaluation}

\subsection{Acceptance Ratio}

Table~\ref{tab:results} presents the comparative evaluation of the Maskable PPO agent against the Best-Fit baseline over 1000 SFC requests. The PPO agent achieves a competitive acceptance ratio, demonstrating that the learned policy is capable of efficiently allocating resources across the substrate network. The Best-Fit baseline, which greedily minimizes co-location and then maximizes security margin at each step, provides a strong comparison point.

\begin{table}[!ht]
\centering
\caption{Comparative Evaluation Results (1000 SFC Requests)}
\label{tab:results}
\begin{tabularx}{\textwidth}{X r r r l}
\hline
\textbf{Metric} & \textbf{Best-Fit mean} & \textbf{PPO mean} & \textbf{$\Delta$\%} & \textbf{Winner} \\ \hline
Acceptance Ratio & 0.9370 & 0.9236 & $-1.43$\% & Best-Fit \\
Avg Security Cost Heuristic & 0.0012 & 0.0010 & $-11.09$\% & PPO \\
Avg Security Margin & 1.0850 & 1.6257 & $+49.84$\% & PPO \\
Total Revenue & 9072.20 & 8965.40 & $-1.18$\% & Best-Fit \\
Avg Revenue / Request & 9.0722 & 8.9654 & $-1.18$\% & Best-Fit \\
Avg SFC Tenancy (SFCs/occupied node) & 1.0226 & 2.2189 & $+116.98$\% & Best-Fit \\
Avg VNF Tenancy (VNFs/occupied node) & 3.0672 & 3.8079 & $+24.15$\% & Best-Fit \\

\hline
\end{tabularx}
\end{table}

\subsection{Summary Bar Chart}

Figure~\ref{fig:summary_bar} provides a visual overview of the mean $\pm$ standard deviation for all eight comparative metrics across evaluation episodes. The error bars indicate episode-to-episode variance, confirming the stability of both algorithms' performance.

\begin{figure}[!ht]
\centering
\includegraphics[width=\textwidth]{compare/summary_comparison.png}
\caption{Grouped bar chart comparing PPO and Best-Fit across all evaluation metrics (mean $\pm$ std over evaluation episodes).}
\label{fig:summary_bar}
\end{figure}

\subsection{Rolling Average Metrics}

Figure~\ref{fig:rolling} presents the rolling average (window 50) of the acceptance ratio for both algorithms over the 1000-request evaluation. The rolling averages reveal the stability of each algorithm's performance over the course of the evaluation and highlight any temporal patterns such as performance degradation under high load.

\begin{figure}[!ht]
\centering
\includegraphics[width=0.8\textwidth]{compare/sfc_ppo_acceptance_ratio.png}
\caption{Rolling acceptance ratio (window = 50) for Maskable PPO and Best-Fit over 1000 requests.}
\label{fig:rolling}
\end{figure}

\subsection{Security Margin}

Figure~\ref{fig:sec_margin} shows the rolling average security margin — the surplus of the placement node's security score over the SFC's minimum requirement — across the evaluation. A higher margin indicates that the algorithm consistently selects nodes with stronger security posture than strictly necessary.

\begin{figure}[!ht]
\centering
\includegraphics[width=0.8\textwidth]{compare/security_score_training.png}
\caption{Rolling average security margin (window = 50) for Maskable PPO and Best-Fit over 1000 requests.}
\label{fig:sec_margin}
\end{figure}

\subsection{Per-Node Tenancy Distribution}

Figures~\ref{fig:tenancy_bar}, \ref{fig:tenancy_hist}, and \ref{fig:tenancy_cdf} show the cumulative per-node SFC and VNF placement counts over the 1000-request evaluation — i.e., how many SFC chains and VNF instances each substrate node hosted in total. This distribution captures which nodes each algorithm tends to prefer and how evenly load is spread across the substrate.

\begin{figure}[!ht]
\centering
\includegraphics[width=\textwidth]{compare/tenancy_grouped_bar.png}
\caption{Cumulative per-node SFC (top) and VNF (bottom) placement counts for Maskable PPO and Best-Fit over 1000 requests. Each bar group represents one substrate node.}
\label{fig:tenancy_bar}
\end{figure}

\begin{figure}[!ht]
\centering
\includegraphics[width=\textwidth]{compare/tenancy_histogram.png}
\caption{Histogram of cumulative per-node SFC (left) and VNF (right) placement counts, showing how many nodes received each count value. A distribution concentrated at high counts indicates hot-spotting; a spread distribution indicates more uniform utilisation.}
\label{fig:tenancy_hist}
\end{figure}

\begin{figure}[!ht]
\centering
\includegraphics[width=\textwidth]{compare/tenancy_cdf.png}
\caption{Empirical CDF of cumulative per-node SFC (left) and VNF (right) placement counts for Maskable PPO and Best-Fit. A CDF that rises steeply at higher count values indicates that a small number of nodes absorb a disproportionate share of placements.}
\label{fig:tenancy_cdf}
\end{figure}

\section{Analysis and Discussion}

\subsection{Benefits of Action Masking}

Action masking provides three key benefits in this framework. First, it guarantees 100\% feasibility of accepted placements: since infeasible nodes are never selected, every placement that completes the chain satisfies all resource, bandwidth, latency, security posture, and isolation constraints. Second, it dramatically improves sample efficiency: without masking, the DRL agent would waste many gradient steps learning to avoid infeasible actions through penalty signals. Third, it simplifies reward design: since constraint violations cannot occur, there is no need for complex penalty terms for individual constraint types.

\subsection{Effect of Security Posture-Aware Reward}

The per-VNF shaping rewards create a soft preference for nodes with high security margin and low co-location, without a hard trade-off against acceptance ratio. The security margin bonus ($\alpha_s = 0.1$) rewards placing on nodes whose effective security exceeds the request minimum, while the co-location penalty ($\alpha_c = 0.1$, inverse saturation) discourages hot-spotting. The windowed AR modifier provides additional stability: when the agent has recently been accepting many requests (high $\overline{\text{AR}}_w$), the acceptance reward is amplified, encouraging continued good performance; when many rejections have occurred (low $\overline{\text{AR}}_w$), the rejection penalty is amplified, pushing the agent to explore better strategies.

\subsection{GNN vs. MLP Policies}

The GNN-based feature extractor captures topology-dependent and security posture-dependent patterns that MLP-based policies cannot represent. For example, the DRL agent may learn to prefer substrate nodes that are central in the graph (with low average distance to all other nodes), reducing latency costs, or nodes with high residual resources, high security scores, and low incident pressure. The attention aggregation mechanism allows the policy to dynamically focus on the most informative nodes — for instance, attending strongly to nodes with degraded security posture when computing the security-aware risk component.

\subsection{Comparison with Best-Fit}

The Best-Fit by Minimum Risk baseline represents a competitive greedy approach that directly minimizes the per-node risk score $\rho_n^{\text{node}}$ at each VNF placement step. The key distinction is that the intelligent DRL agent optimizes security posture \textit{jointly with} acceptance ratio and long-term resource efficiency via the shaping rewards, while the baseline myopically minimizes risk at each individual step. In heavily loaded scenarios, the DRL agent may make suboptimal individual-step choices early in the episode to preserve high-security, high-resource nodes for future requests with stricter constraints — a strategy the greedy baseline cannot exhibit.

\textbf{Tenancy spread.}
The per-node tenancy distribution (Figures~\ref{fig:tenancy_bar}--\ref{fig:tenancy_cdf}) reveals that PPO concentrates placements on a smaller set of preferred nodes compared to Best-Fit, exhibiting a \textit{hot-spotting} pattern. The co-location shaping penalty uses an inverse saturation $1 - 1/(1 + \text{SFC}(n))$ which increases steeply for the first few chains but flattens quickly — a node with $\text{SFC}(n) = 5$ receives only marginally more penalty than one with $\text{SFC}(n) = 3$. On high-security nodes the security margin bonus simultaneously provides a persistent pull. Best-Fit, by contrast, selects $\arg\min_n \rho_n^{\text{node}} = q_n \cdot \text{SFC}(n) \cdot T_s/\tau_{\max}$ at each step: as $\text{SFC}(n)$ accumulates on any node, that node's score grows linearly and previously under-used nodes become comparatively more attractive, producing a natural anti-concentration effect. Concentration on a small subset of nodes increases blast radius in the event of a correlated failure; a stricter per-node SFC hard limit in the reward would be a straightforward countermeasure.

\subsection{Scalability Considerations}

The current implementation uses a 25-node substrate, which enables efficient training within 100,000 timesteps. The GNN architecture is topology-agnostic by design: the `node\_mask` ensures that only real nodes are processed, and the edge index is reconstructed dynamically from the substrate graph. This means the same trained DRL policy can, in principle, be deployed on larger or reconfigured substrates, maintaining its security posture awareness across different network configurations.
